\section{Results}\label{results}

\subsection{Predictive accuracy}\label{predictive-accuracy}

The reproducible current fixed RF on AE-dB-z + Vib-dB-z gives a mean
LOGO MAE of \(0.0210~\upmu\)m, median \(0.0093~\upmu\)m, and maximum
\(0.1478~\upmu\)m. Cache-matched inner-grouped runs give 0.0207~\(\upmu\)m
for dB-z and 0.0206~\(\upmu\)m for log-mel, supporting a leading
representation group by mean MAE rather than a unique winner. Their tails are
not equivalent: the current nested maximum fold MAE is 0.1402~\(\upmu\)m for
dB-z and 0.2061~\(\upmu\)m for log-mel. Thus, dB-z has materially better
observed worst-condition robustness in this 16-fold analysis. Table~\ref{tbl:top10}
and Figure~\ref{fig:ranking} retain the historical canonical benchmark for
comparison with the archived 14-test family; those artifacts cannot be
regenerated exactly in the current environment. In that historical ranking,
the dB-z RF gives 0.0200~\(\upmu\)m, median 0.0089~\(\upmu\)m, and maximum
0.1437~\(\upmu\)m. In both current and historical pipelines the mean is dominated by a single
difficult condition. The
largest single-fold errors are driven by condition 7 (see
Section~\ref{sec:conditions} and Figure~\ref{fig:top-models-fold-mae}).
Configuration abbreviations are defined in
Table~\ref{tbl:config-abbreviations}.

\begin{table}[htbp]
\centering
\caption{Primary reproducible current RF results. Fixed and nested runs use
the same 14 training conditions in each outer fold; nested runs differ only
through inner grouped hyperparameter selection.}
\label{tbl:current-rf-primary}
\begin{tabular}{@{}lcccc@{}}
\toprule
Representation/protocol & Training conditions & Mean & Median & Maximum \\
\midrule
dB-z, fixed RF & 14 & 0.0210 & 0.0093 & 0.1478 \\
dB-z, inner-grouped RF & 14 & 0.0207 & 0.0093 & 0.1402 \\
Log-mel, inner-grouped RF & 14 & 0.0206 & 0.0087 & 0.2061 \\
\bottomrule
\end{tabular}
\end{table}

\begin{figure}[htbp]
\centering
\includegraphics[width=183mm]{nested_rf_dbz_logmel_folds.png}
\caption{Paired outer-fold MAE for the current nested dB-z and implemented
log-mel RF analyses. Their means are nearly identical, but log-mel has the
larger Condition~7 and maximum-fold error.}\label{fig:nested-rf-tail}
\end{figure}

Table~\ref{tbl:top10} should be read as a post-selection ranking because
model family and representation vary together. First, multimodal fusion gives the best point
estimate, but the gain over the best vibration-only configuration
(Random Forest on Vib-log-mel, rank~7; MAE 0.0211~\(\upmu\)m, and the
vibration-only ablation in Section~\ref{sec:robustness}, MAE
0.0207~\(\upmu\)m) is small and should not be interpreted as a
statistically established multimodal advantage. Second, per-sample-standardised full-resolution dB descriptors and
log-mel descriptors occupy the top ranks, but the difference
relative to dB spectrograms is small
and not statistically significant under the present LOGO protocol (see
Section~\ref{statistical-comparison}). Third, adding process parameters
(PP) to the top dB-z/log-mel configurations has a negligible effect (MAE
changes by \(<0.0006~\upmu\)m). Because each condition is a unique
process-parameter combination, LOGO tests generalization to withheld
process-parameter combinations within the sampled factor space; the poor
added value of PP may therefore reflect the sparse 16-condition design
rather than true physical irrelevance.

\subsubsection{Current fixed-setting repeated-baseline sensitivity of stochastic
architectures}\label{sec:seed-repeat}

To assess stochastic variability under a common current protocol, the main stochastic architectures
were each repeated over five random seeds under the same 16 LOGO folds
(Supplementary Table~\ref{tab:supp-seed-level}). These runs do not exactly
reproduce the historical canonical checkpoints and should therefore be
interpreted as a current fixed-setting stochastic sensitivity, not as variance
around the historical tuned seed-42 models. The repeated neural models
and RF use the same canonical 14-condition training budget in each fold. For
the current random forest on AE-dB-z + Vib-dB-z, the seed-averaged MAE is
\(0.0203 \pm 0.0360~\upmu\)m, with per-seed overall MAEs of 0.0210,
0.0205, 0.0204, 0.0201, and 0.0197~\(\upmu\)m. The seed-level standard
deviation of these five overall MAEs is 0.00043~\(\upmu\)m, showing
that the RF result is stable across seeds; this internal stability does
not, however, address measurement uncertainty or external generalisation. The deep-learning baselines
exhibit larger seed-to-seed instability: ResNetVibCNN seed-level MAEs
range from 0.0283 to 0.0441~\(\upmu\)m, ResNetFusion from 0.0369 to
0.0487~\(\upmu\)m, BilinearFusionNetwork from 0.0414 to 0.0492~\(\upmu\)m,
TrajectoryCNN from 0.0439 to 0.0484~\(\upmu\)m, ResNetAECNN from 0.0353 to
0.0478~\(\upmu\)m, and ChannelAttentionCNN from 0.0334 to 0.0503~\(\upmu\)m.
Their seed-averaged MAEs (ResNetVibCNN 0.0384, ResNetFusion 0.0417,
BilinearFusionNetwork 0.0447, TrajectoryCNN 0.0450, ResNetAECNN 0.0423,
and ChannelAttentionCNN 0.0423~\(\upmu\)m) differ from the historical
canonical point estimates. Within the repeated set, seed 42 is not uniformly
the lowest-error seed. Supplementary
Table~\ref{tab:supp-seed-repeat-all} gives the per-condition
seed-averaged errors for the equal-budget current random forest and the six
repeated deep-learning baselines. The seed-averaged per-condition
MAE for the current random forest remains dominated by Condition~7
(0.1509~\(\upmu\)m), followed by Condition~8 (0.0585~\(\upmu\)m), so the
seed-repeat analysis shares the same failure mode as the primary
benchmark. The $\pm$ values in the seed-averaged MAE paragraph refer to
fold-level variability across the 16 LOGO conditions; seed-level
variability is reported separately through the per-seed overall MAEs.

Paired Wilcoxon signed-rank tests on the 16 condition-level seed-averaged
MAEs show that the random forest is significantly lower than the six
repeated deep-learning baselines after Holm--Bonferroni correction for the
six comparisons (all \(p_{\text{Holm}} < 0.01\); the largest corrected
\(p\)-value is 0.0027 versus TrajectoryCNN). This equal-budget sensitivity analysis
provides robustness evidence within the selected current fixed-setting
baselines. Because the historical tuned checkpoints are not reproduced, the
differences reflect the current protocol as well as seed variation and do not
independently confirm superiority over those historical neural models. The inference
family here contains only the six selected repeated stochastic baselines
and should not be combined with the 14-comparison primary family. Because
the repeated model set was selected from the broader benchmark, this
result is reported as robustness evidence rather than as the primary
inferential comparison.

The validation-condition rotation diagnostic provides a separate check on
the deep early-stopping split. With fixed default architectures and training
settings, ResNetVibCNN changes from 0.0393~\(\upmu\)m under the canonical
seed-42 validation draw to 0.0471~\(\upmu\)m under the cyclic-next
assignment; ResNetAECNN is similar at 0.0440 and 0.0435~\(\upmu\)m,
respectively (Supplementary Table~\ref{tab:supp-deep-validation-rotation}).
These values are a validation-assignment sensitivity analysis, not a
replacement for the separately tuned canonical benchmarks, because the
original per-fold tuning states are not reconstructed here.

Table~\ref{tbl:top10-summary} supplements the mean and standard
deviation with median MAE, interquartile range, bootstrap 95\% CI, and
the minimum and maximum fold MAE for the top five models. The bootstrap
CI is computed by resampling the 16 LOGO folds with replacement. The
summary confirms that the top models are strongly tail-sensitive:
the maximum fold MAE is an order of magnitude larger than the median for
every leading model. Because the expanded profilometer uncertainty is
0.025~\(\upmu\)m, the best MAE of 0.020~\(\upmu\)m should be interpreted as
error against the averaged measured \(R_a\), not against latent true
roughness. Raw trace-level repeatability is unavailable, so the reported
MAE combines model error and measurement uncertainty; the true
roughness-prediction error cannot be separated from measurement
uncertainty, and ranking differences smaller than approximately
0.005--0.010~\(\upmu\)m should be treated as exploratory.

\begin{table}[htbp]
\centering
\caption{Historical canonical benchmark: distribution of fold-level MAE for the top five
model--configuration pairs. These archived results are not the reproducible
current RF reruns. The 10\% trimmed mean removes the lowest and
highest folds and is descriptive only; CI = bootstrap 95\% confidence
interval for the mean obtained by resampling folds.}
\label{tbl:top10-summary}
\begin{tabular}{@{}lccccccc@{}}
\toprule
Model / Config & Mean & Median & IQR & 10\% trim. & Min & Max & 95\% CI \\
\midrule
Random forest / AE-dB-z + Vib-dB-z & 0.0200 & 0.0089 & 0.0055 & 0.0124 & 0.0031 & 0.1437 & 0.0077--0.0405 \\
Random forest / AE-log-mel + Vib-log-mel & 0.0201 & 0.0084 & 0.0044 & 0.0083 & 0.0048 & 0.2010 & 0.0072--0.0448 \\
Random forest / AE-dB-z + Vib-dB-z + PP & 0.0205 & 0.0091 & 0.0047 & 0.0127 & 0.0032 & 0.1472 & 0.0076--0.0417 \\
Random forest / AE-log-mel + Vib-log-mel + PP & 0.0205 & 0.0091 & 0.0053 & 0.0084 & 0.0046 & 0.2067 & 0.0072--0.0459 \\
Random forest / AE-dB + Vib-dB + PP & 0.0210 & 0.0094 & 0.0039 & 0.0110 & 0.0033 & 0.1785 & 0.0084--0.0437 \\
\bottomrule
\end{tabular}
\end{table}

The ranking should be interpreted using both mean and worst-condition
performance because the mean is dominated by Condition~7. The distribution
is strongly tail-sensitive: the median fold MAE for the best RF is
0.0089~\(\upmu\)m, but the maximum fold MAE is 0.1437~\(\upmu\)m
(Condition~7), so the mean is dominated by a single difficult condition. The 10\% trimmed mean (excluding the lowest and
highest folds) is 0.0124~\(\upmu\)m, and the 20\% trimmed mean is
0.0090~\(\upmu\)m, confirming that typical-condition performance is
lower-error than the mean, while worst-condition performance remains the
practical limiting case. These trimmed means are descriptive only and do not
replace the main LOGO mean.

Figure~\ref{fig:ranking} extends this historical view to the top 20. Auditable
tree-ensemble model--configuration pairs occupy the top ranks, with the
first nine places held exclusively by random forests. The first non-tree
model, Trajectory CNN on Vib-traj, appears at rank 13 (MAE
\(= 0.0259 \pm 0.0565~\upmu\)m), about 30\% higher in point estimate
than the best random forest, although the primary corrected paired
comparison is not statistically significant. The best spectrogram-based
deep CNN, ResNetVibCNN on Vib-dB, is at rank 14 (MAE
\(= 0.0268 \pm 0.0443~\upmu\)m), about 34\% higher in point estimate
and likewise non-significant after Holm--Bonferroni correction.
This ordering is noteworthy because the deep models have tens of
thousands of gradient-trained weights, whereas the tree ensembles store
split thresholds and leaf statistics rather than conventional network
weights. The ranking should not, however, be read as a controlled
model-only comparison: model family and representation are partly
confounded, because a model may rank highly in part because it was paired
with a stronger representation.

A model-selection caveat applies to this comparison. The deep-learning
entries were tuned by random search over model and training hyperparameters
on the same outer LOGO folds, while the tree-based models used fixed
hyperparameters (Table~\ref{tbl:hyperparams}). This non-nested design may
favour the deep-learning estimates \cite{varma2006bias}; the ranking should
therefore be interpreted as a relative performance ordering under a fixed
evaluation budget rather than an unbiased estimate of population
generalisation.

\begin{table}[htbp]
\centering
\caption{Historical canonical benchmark: top 10 model--configuration pairs
ranked by LOGO MAE. Current reproducible fixed and inner-grouped RF results
are reported in the opening paragraph of Section~\ref{predictive-accuracy}.}
\label{tbl:top10}
\begin{tabular}{@{}rlcc@{}}
\toprule
Rank & Model & Configuration & MAE (\(\upmu\)m) \\
\midrule
1 & Random forest & AE-dB-z + Vib-dB-z & 0.0200 \\
2 & Random forest & AE-log-mel + Vib-log-mel & 0.0201 \\
3 & Random forest & AE-dB-z + Vib-dB-z + PP & 0.0205 \\
4 & Random forest & AE-log-mel + Vib-log-mel + PP & 0.0205 \\
5 & Random forest & AE-dB + Vib-dB + PP & 0.0210 \\
6 & Random forest & AE-dB + Vib-dB & 0.0210 \\
7 & Random forest & Vib-log-mel & 0.0211 \\
8 & Random forest & Vib-dB & 0.0221 \\
9 & Random forest & Vib-dB-z & 0.0223 \\
10 & Ridge regression & AE-dB & 0.0246 \\
\bottomrule
\end{tabular}
\end{table}

Because the best reported MAE (0.020~\(\upmu\)m) is only about twice the
manufacturer-quoted profilometer repeatability and raw trace-level
repeatability is unavailable, differences smaller than approximately
0.005--0.010~\(\upmu\)m among the top RF variants should be treated as
exploratory.

\begin{figure}
\centering
\includegraphics[width=183mm]{submission_ranking_top20.png}
\caption{Historical canonical top-20 model-configuration pairs ranked by LOGO MAE.
Auditable tree-ensemble model--configuration pairs occupy the upper
ranks by point-estimate MAE, although model and representation effects
are partly confounded.}\label{fig:ranking}
\end{figure}

Figure~\ref{fig:scatter-top3} shows predicted versus observed surface
roughness for the three best-performing configurations. The points cluster
around the ideal \(y=x\) line across most of the measured range,
consistent with useful predictive accuracy. The scatter is wider at
higher roughness values, where condition 7 contributes several large
residuals (see Section~\ref{sec:conditions}). Residuals are generally
small within the training-target support, but they show marked
regime-dependent bias outside that support. Condition~7 is
systematically underpredicted, consistent with regression toward the
training target range. Overall, the random forest on AE-dB-z +
Vib-dB-z is one of the leading point-estimate candidates, but its fold-level
variance limits how strongly its condition-wise reliability can be
asserted.

\begin{figure}
\centering
\includegraphics[width=183mm]{submission_prediction_scatter_top3.png}
\caption{Historical canonical benchmark predictions versus observed surface
roughness for the top three model-configuration pairs. The dashed black line marks ideal
\(y=x\) calibration.}\label{fig:scatter-top3}
\end{figure}

To test whether these accuracy differences are statistically robust, the
next subsection reports a systematic pairwise comparison of the leading
models.

\subsection{Primary post-selection statistical comparison}\label{statistical-comparison}

The planned family used for the historical canonical analysis contains 14
Holm--Bonferroni-corrected Wilcoxon signed-rank tests on
per-fold LOGO MAEs (Section~\ref{validation-protocol}); the five-seed
sensitivity analysis in Section~\ref{sec:seed-repeat} is reported
separately as robustness evidence and is not included in this family.
Because model and representation selection used observed LOGO results, the
paired tests are post-selection diagnostics rather than confirmatory
population-level inference. To assess whether the accuracy
differences in Section~\ref{predictive-accuracy} are robust, we performed
nine of the planned pairwise Wilcoxon signed-rank tests on per-fold LOGO
MAEs; the raw \(p\)-values were corrected with Holm--Bonferroni across
the full family of 14 planned comparisons (Table~\ref{tbl:statistical};
the complete family is given in Supplementary
Table~\ref{tab:supp-holm}). The nine comparisons shown here isolate
specific effects: the strongest auditable classical model versus the strongest
deep CNN; the best random forest versus the best wavelet-scattering
model; the same fused spectrogram input given to a random forest or a
ridge regressor; a random forest on spectrograms versus a shallow MLP on
hand-crafted features; the same vibration spectrogram given to a random
forest or ResNetVibCNN; and dB-z versus dB spectrograms within the
random forest.

\subsubsection{Pairwise model comparisons}\label{sec:pairwise-comparisons}

\begin{table}[htbp]
\centering
\caption{Historical canonical benchmark pairwise statistical tests (Holm-Bonferroni corrected
across the full family of 14 planned comparisons). HL \(\Delta\) is the
Hodges--Lehmann paired location shift for model~A minus model~B, with an
unadjusted descriptive bootstrap 95\% confidence interval; \(r_{rb}\) is the Wilcoxon
rank-biserial correlation. Negative values favour model~A.}\label{tbl:statistical}
\resizebox{\linewidth}{!}{\begin{tabular}{@{}lccccccc@{}}
\toprule
Comparison & MAE A & MAE B & HL $\Delta$ & 95\% CI & $r_{rb}$ & $p$ raw & $p$ Holm \\
\midrule
Best RF (auditable ML) vs Best DL & 0.0200 & 0.0268 & -0.0042 & [-0.0164, 0.0003] & -0.49 & 0.0934 & 1.0000 \\
Best RF (spectrogram) vs Best WST & 0.0200 & 0.0225 & -0.0011 & [-0.0039, 0.0018] & -0.26 & 0.3755 & 1.0000 \\
RF dB vs Ridge dB & 0.0210 & 0.0248 & -0.0050 & [-0.0097, -0.0012] & -0.60 & 0.0335 & 0.4360 \\
RF dB vs Shallow MLP features & 0.0210 & 0.0495 & -0.0249 & [-0.0412, -0.0127] & -1.00 & 3.05e-5 & 4.27e-4 \\
RF Vib-dB vs ResNetVibCNN Vib-dB & 0.0221 & 0.0268 & -0.0032 & [-0.0086, -0.0001] & -0.54 & 0.0577 & 0.6921 \\
RF dB-z vs RF dB & 0.0200 & 0.0210 & -0.0006 & [-0.0031, 0.0017] & -0.18 & 0.5619 & 1.0000 \\
RF AE-dB vs ResNetAECNN & 0.0403 & 0.0300 & 0.0045 & [-0.0016, 0.0266] & 0.34 & 0.2522 & 1.0000 \\
ResNetAECNN vs ChannelAttentionCNN & 0.0300 & 0.0406 & -0.0098 & [-0.0229, 0.0029] & -0.37 & 0.2114 & 1.0000 \\
ResNetVibCNN vs ResNetAECNN & 0.0268 & 0.0300 & -0.0002 & [-0.0123, 0.0093] & -0.01 & 0.9799 & 1.0000 \\
\bottomrule
\end{tabular}
}
\end{table}

Only one comparison survives Holm-Bonferroni correction: the random
forest on AE-dB + Vib-dB is markedly better than the shallow MLP on
AE/Vib feats + PP (Hodges--Lehmann paired difference
\(-0.0249~\upmu\)m, \(p_{\text{Holm}}=4.27\times10^{-4}\)). This is a
comparison of the archived implementations: the historical MLP used an
internal random validation fraction rather than the designated condition.
The primary statistical family is therefore an artifact-based post-selection
analysis of preserved historical canonical predictions, not an independently
regenerated comparison.
A corrected grouped-validation, target-scaled shallow MLP gives mean MAE
0.0529~\(\upmu\)m (maximum 0.2645~\(\upmu\)m), compared with 0.0210 for
the current RF; the descriptive paired Wilcoxon raw \(p\)-value is
\(3.05\times10^{-5}\) (Supplementary
Table~\ref{tab:supp-grouped-mlp}). This sensitivity supports the observed
separation but is not part of the historical Holm family. All other
displayed comparisons are non-significant after correction. Most favour
the auditable or simpler model by point estimate, but not all:
RF on AE-dB versus ResNetAECNN favours the CNN by point estimate
(\(p_{\text{Holm}}=1.0\)). The signed-rank tests use
only 16 paired observations, so the study is underpowered for small
model differences; non-significant results should not be interpreted as
equivalence. The reported Hodges--Lehmann confidence intervals and
rank-biserial correlations are standard non-parametric effect-size
summaries, but the effective paired sample size is 16 conditions and the
confidence intervals quantify fold-level uncertainty within this design,
not population-level uncertainty. The Holm-Bonferroni
adjustment was applied across the full family of 14 planned
comparisons (Table~\ref{tbl:statistical} and Supplementary
Table~\ref{tab:supp-holm}); the corrected values should be interpreted
within this family. The lack of significance is driven by high fold-level
variance: condition 7 (Section~\ref{sec:conditions}) inflates the MAE of
every model and is the largest contributor to the LOGO error budget.

The closest comparison is the random forest versus ResNetVibCNN on the
same Vib-dB input (HL \(\Delta=-0.0032~\upmu\)m,
\(p_{\text{Holm}}=0.6921\)). Although the point estimate favours the
random forest, the wide fold-level variance prevents a definitive
conclusion. The same pattern holds for the best random forest versus the
best spectrogram-based CNN (HL \(\Delta=-0.0042~\upmu\)m,
\(p_{\text{Holm}}=1.0000\)) and versus the best gradient-boosted model
(LightGBM on Vib-dB-z, reported in Supplementary
Table~\ref{tab:supp-holm}, HL \(\Delta=-0.0008~\upmu\)m,
\(p_{\text{Holm}}=1.0\)). Even the
ridge regressor on the same fused spectrogram is not significantly worse
than the random forest (\(p_{\text{Holm}}=0.4360\)), suggesting that the
representation itself carries most of the predictive signal. The
dB-z versus dB-spectrogram comparison within the random forest is
essentially a null result (HL \(\Delta=-0.0006~\upmu\)m,
\(p_{\text{Holm}}=1.0\)); the small ranking advantage of the dB-z cache
seen in Table~\ref{tbl:top10} is therefore not statistically decisive
under LOGO.

Because condition 7 is the dominant source of fold-level variance,
Table~\ref{tbl:condition7-sensitivity} reports the median paired
difference, Wilcoxon \(p\)-value, and bootstrap 95\% confidence interval
for three key comparisons with and without condition 7. Excluding
condition 7 narrows the point estimates but does not change the
qualitative conclusion: the random forest remains favourable by point
estimate, yet the reduction in variance is modest because other
conditions also contribute spread. This sensitivity analysis is
illustrative only; condition 7 is retained in all main reported metrics.

\begin{table}[htbp]
\centering
\caption{Exploratory sensitivity of paired differences to excluding
condition 7. Negative \(\Delta\) favours model A. Values are median
paired LOGO MAE differences (\(\upmu\)m) with unadjusted Wilcoxon
\(p\)-values and bootstrap 95\% CIs; these \(p\)-values are reported only
for post-hoc sensitivity analysis.}
\label{tbl:condition7-sensitivity}
\begin{tabular}{@{}lcccccc@{}}
\toprule
& \multicolumn{3}{c}{All folds} & \multicolumn{3}{c}{Excl. condition 7} \\
\cmidrule(lr){2-4} \cmidrule(lr){5-7}
Comparison & \(\Delta\) & \(p\) & 95\% CI & \(\Delta\) & \(p\) & 95\% CI \\
\midrule
Best RF vs RF AE-dB + Vib-dB + PP & \(-0.0005\) & 0.528 & \(-0.0024, 0.0011\) & \(-0.0002\) & 0.804 & \(-0.0014, 0.0031\) \\
Best RF vs ResNetVibCNN (Vib-dB) & \(-0.0042\) & 0.093 & \(-0.0083, 0.0015\) & \(-0.0040\) & 0.169 & \(-0.0071, 0.0029\) \\
Best RF vs LightGBM (Vib-dB-z) & \(-0.0005\) & 0.375 & \(-0.0026, 0.0007\) & \(-0.0002\) & 0.599 & \(-0.0022, 0.0007\) \\
\bottomrule
\end{tabular}
\end{table}

\begin{figure}
\centering
\includegraphics[width=183mm,height=\textheight,keepaspectratio]{results_statistical_comparison_pairs.png}
\caption{Historical canonical benchmark: paired mean LOGO MAE and fold
standard deviation for the nine statistical comparisons in
Table~\ref{tbl:statistical}. Comparison IDs
C1--C9 follow the table order. Stars mark Holm-Bonferroni significance
(*~\(p<0.05\)).}\label{fig:statistical-pairs}
\end{figure}

Figure~\ref{fig:statistical-pairs} restates the nine displayed comparisons as
paired mean MAEs with fold-level standard deviations. Most auditable classical or
simpler configurations sit on the lower-error side of their pair, but the
AE-only RF versus ResNetAECNN comparison is an exception. The overlapping
error bars explain why only the shallow-MLP comparison is formally
significant.

\begin{figure}
\centering
\includegraphics[width=183mm,height=\textheight,keepaspectratio]{results_statistical_comparison_forest.png}
\caption{Historical canonical benchmark forest plot of median paired LOGO MAE differences for the nine
comparisons in Table~\ref{tbl:statistical}, with descriptive 95\% bootstrap
confidence intervals. These intervals are unadjusted and should not be interpreted as Holm--Bonferroni significance. Negative values favour Model~A.}\label{fig:statistical-forest}
\end{figure}

Figure~\ref{fig:statistical-forest} recasts the same comparisons as a
forest plot of median paired differences with descriptive 95\% bootstrap confidence
intervals. Two unadjusted descriptive intervals exclude zero: RF dB spectrogram versus Ridge dB spectrogram and RF dB spectrogram versus shallow MLP. Only the RF-versus-MLP comparison remains significant after Holm--Bonferroni correction; the Holm-corrected inference comes from Table~\ref{tbl:statistical}, rather than from the plotted intervals. Taken together, the
tests support clear superiority over the shallow MLP baseline, while the
remaining comparisons show favourable but non-significant point-estimate
differences. Because the remaining variance appears to be driven by how
sensor information is represented, the next subsection compares the
signal-representation families directly.

\subsection{Signal-representation
comparison}\label{signal-representation-comparison}

Table~\ref{tbl:representations} summarises the historical canonical best LOGO-only result
obtained with each signal-representation family, and
Figure~\ref{fig:representations-bar} visualises the corresponding mean
LOGO MAEs. Because the best-performing model within each family was
selected separately, representation and model choice are confounded in
this table; it shows the best observed model--representation combination
within each family rather than a controlled representation-only
comparison. Table~\ref{tbl:controlled-representations} and
Figure~\ref{fig:representation-controlled} instead hold the exact
input configuration fixed within each representation row for Random
Forest, Ridge, and LightGBM. They are a controlled
learner-by-representation sensitivity analysis, not a post-selection
ranking. The observed values span a substantial
range across the three stable learners. The fixed-input rows are a
representation sensitivity analysis, not a ranking that identifies a
unique best representation. Preprocessing parameters for these
representations are given in Section~\ref{signal-representations}.

\begin{table}[htbp]
\centering
\caption{Historical canonical best-observed LOGO-only result per signal
representation. These values are retained for benchmark provenance rather
than as the reproducible current representation ranking.}\label{tbl:representations}
\begin{tabular}[]{@{}
  >{\raggedright\arraybackslash}p{(\linewidth - 10\tabcolsep) * \real{0.2000}}
  >{\raggedright\arraybackslash}p{(\linewidth - 10\tabcolsep) * \real{0.1800}}
  >{\raggedright\arraybackslash}p{(\linewidth - 10\tabcolsep) * \real{0.1600}}
  >{\raggedright\arraybackslash}p{(\linewidth - 10\tabcolsep) * \real{0.1600}}
  >{\raggedright\arraybackslash}p{(\linewidth - 10\tabcolsep) * \real{0.1700}}
  >{\raggedright\arraybackslash}p{(\linewidth - 10\tabcolsep) * \real{0.0800}}@{}}
\toprule
\begin{minipage}[b]{\linewidth}\raggedright
Representation
\end{minipage} & \begin{minipage}[b]{\linewidth}\raggedright
Model
\end{minipage} & \begin{minipage}[b]{\linewidth}\raggedright
Config
\end{minipage} & \begin{minipage}[b]{\linewidth}\raggedright
MAE (\(\upmu\)m)
\end{minipage} & \begin{minipage}[b]{\linewidth}\raggedright
MSE (\(\upmu\)m\(^2\))
\end{minipage} & \begin{minipage}[b]{\linewidth}\raggedright
Folds
\end{minipage} \\
\midrule
dB-z spectrogram & Random forest & \texttt{AE-dB-z + Vib-dB-z}
& \(0.0200 \pm 0.0356\) & \(0.0016 \pm 0.0052\) & 16 \\
Log-mel spectrogram & Random forest & \texttt{AE-log-mel + Vib-log-mel} &
\(0.0201 \pm 0.0483\) & \(0.0028 \pm 0.0107\) & 16 \\
Spectrogram & Random forest & AE-dB + Vib-dB + PP &
\(0.0210 \pm 0.0425\) & \(0.0023 \pm 0.0083\) & 16 \\
Wavelet scattering & LightGBM & \texttt{Vib-WST} &
\(0.0225 \pm 0.0358\) & \(0.0019 \pm 0.0050\) & 16 \\
Trajectory & Trajectory CNN & \texttt{Vib-traj} &
\(0.0259 \pm 0.0565\) & \(0.0040 \pm 0.0140\) & 16 \\
Time-domain features & Bilinear fusion &
\texttt{AE/Vib feats + phys + PP} & \(0.0358 \pm 0.0424\) &
\(0.0034 \pm 0.0073\) & 16 \\
Process parameters & TabTransformer & \texttt{Process params} &
\(0.0418 \pm 0.0672\) & \(0.0060 \pm 0.0162\) & 16 \\
\bottomrule
\end{tabular}
\end{table}

The exact input configurations in
Table~\ref{tbl:controlled-representations} are, by row, AE-dB + Vib-dB,
AE-dB-z + Vib-dB-z, AE-log-mel + Vib-log-mel, AE-WST + Vib-WST,
AE/Vib time-domain features, and process parameters. The same 16 seed-42
LOGO folds, auxiliary-input policy, and fixed learner hyperparameters are
used for all cells. This design isolates the representation/input choice
within each learner more directly than the best-observed ranking in
Table~\ref{tbl:representations}; it does not turn the broader benchmark
into nested confirmatory inference.

\begin{figure}
\centering
\includegraphics[width=183mm]{methods_signal_representation_comparison.png}
\caption{Historical canonical best-observed mean LOGO MAE for each signal representation family.
Error bars show fold standard deviation across the 16 LOGO folds; lower
whiskers are capped at zero because MAE is non-negative.}
\label{fig:representations-bar}
\end{figure}

\begin{table}[htbp]
\centering
\caption{Fixed-input representation comparison across three stable tabular
learners. Every cell within a row uses the exact configuration named in
the first column, the same 16 seed-42 LOGO folds, and no auxiliary inputs
unless they are named. Values are mean LOGO MAE across folds; fold standard
deviations are shown in Figure~\ref{fig:representation-controlled}. The
shallow MLP is excluded because its raw-target fixed-input fits failed
numerically; a target-scaled and tuned MLP comparison would be a separate
experiment.}
\label{tbl:controlled-representations}
\begin{tabular}{@{}lccc@{}}
\toprule
Fixed input & RF & Ridge & LightGBM \\
\midrule
AE-dB + Vib-dB & 0.0212 & 0.0248 & 0.0321 \\
AE-dB-z + Vib-dB-z & 0.0210 & 0.0237 & 0.0326 \\
AE-log-mel + Vib-log-mel & 0.0205 & 0.0396 & 0.0287 \\
AE-WST + Vib-WST & 0.0331 & 0.0329 & 0.0302 \\
AE/Vib time-domain features & 0.0493 & 0.0451 & 0.0453 \\
Process parameters & 0.0803 & 0.0591 & 0.0733 \\
\bottomrule
\end{tabular}
\end{table}

\begin{figure}
\centering
\includegraphics[width=183mm,height=\textheight,keepaspectratio]{representation_controlled_comparison.png}
\caption{Fixed-input mean LOGO MAE comparison across three stable tabular
learners. Each row uses the exact AE/vibration or process-parameter input
named in Table~\ref{tbl:controlled-representations}; cells show mean
MAE $\pm$ fold standard deviation across the 16 seed-42 LOGO folds.
Trajectory is not shown because no common trajectory input exists for all
three fixed learners.}
\label{fig:representation-controlled}
\end{figure}

The shallow MLP is excluded from this fixed-input table because its
raw-target high-dimensional fits failed numerically. Its values would not
be evidence about representation quality; a target-scaled and tuned MLP
comparison is required before it can be reintroduced. The fixed-input
comparison is therefore most informative for the representation sensitivity
of the three stable tabular learners.

\subsubsection{Modality and fusion effects}\label{sec:modality-fusion}

The leading canonical point estimates are pass-level fused AE and vibration
dB-z descriptors (0.0200~\(\upmu\)m) and log-mel descriptors
(0.0201~\(\upmu\)m), followed by fused dB descriptors with process
parameters (0.0210~\(\upmu\)m). The documented current dB-z pipeline
regenerates 0.0210~\(\upmu\)m rather than the canonical artifact, so this
near-tie should be treated as an observed top group rather than a unique
representation winner. A current-cache inner-grouped sensitivity gives
0.0207~\(\upmu\)m for dB-z and 0.0206~\(\upmu\)m for log-mel RF inputs
(Supplementary Table~\ref{tab:supp-nested-logo}). After inner selection,
each RF is refitted on the same 14 conditions used by the fixed RF; the
archived validation condition remains excluded, so this comparison does not
gain a fifteenth training condition. This independently
supports a leading representation group rather than a unique winner.
This clustering suggests that frequency-domain content carries much
of the predictive information, although model and representation effects
are not fully separated. The vibration-only wavelet-scattering transform gives a close
point estimate (0.0225~\(\upmu\)m), although the difference relative to
pass-level fused dB-z descriptors is not statistically decisive. The trajectory CNN (0.0259~\(\upmu\)m) sits in a middle tier:
it outperforms hand-crafted features and process parameters but does
not reach the spectrogram family. Under the evaluated architectures and
training budget, the trajectory configuration performs less well than the
leading spectral configurations.

Time-domain statistics (0.0358~\(\upmu\)m) and process parameters alone
(0.0418~\(\upmu\)m) fall substantially behind. Their lower performance is
consistent with loss of frequency-localised information, although feature
construction, learner choice, dimensionality, and tuning are not fully
separated. The evaluated scalar statistics compress the selected
steady-state signal windows and may therefore discard localised spectral
structure, while process parameters only capture the nominal grinding state
and cannot resolve condition-specific dynamics. Taken together, these results suggest that,
among the evaluated configurations, frequency-domain vibration-dominant
inputs provide the strongest signal; adding AE gives only a small
mean-level improvement in the recommended RF configuration.

\subsubsection{Time--frequency representation and preprocessing effects}\label{sec:preprocessing-effects}

The observed ranking is consistent with several preprocessing effects,
although these are not isolated experimentally. A controlled comparison
of STFT magnitude, STFT power, log-STFT, mel, and log-mel variants would
be needed to determine whether dynamic-range compression or frequency
binning drives the RF advantage. The small gap between dB spectrograms and
wavelet scattering (0.0210 versus 0.0225~\(\upmu\)m) shows that
frequency-domain inputs are consistently competitive, but the relative
contributions of frequency resolution, compression, normalisation, and
learner choice remain unresolved. A strict ablation removing per-sample/channel standardisation from the
top RF/dB-z configuration leaves the overall LOGO MAE essentially
close (0.0210~\(\upmu\)m with per-sample standardisation and
0.0213~\(\upmu\)m without it), so frequency-shape
information appears sufficient for the mean MAE of this RF/dB-z
configuration. The worst-condition tail is not unchanged, however: the
maximum condition MAE increases from 0.1478 to 0.1863~\(\upmu\)m
(Supplementary Table~\ref{tab:supp-psnorm-ablation}). The ablation therefore supports
mean-level robustness but not tail robustness, and it was not
performed for all models, representations, or uncertainty analyses, so the
conclusion is limited to this RF configuration. Third, adding process parameters or physics
features to the top dB-z/log-mel configurations has a small effect (MAE
changes by \(<0.0006~\upmu\)m), suggesting that the spectrograms already
encode most of the condition-related information in this sparse design.
This result should not be interpreted as physical irrelevance of process
parameters; it may reflect the sparse fractional design and the fact
that condition information is indirectly encoded in the sensor signals.
Because frequency-domain content appears to be the strongest source of
predictive information among the evaluated configurations, the next
question is whether the models exploit physically plausible frequency
regions.

\subsection{Interpretability and physics
consistency}\label{interpretability-and-physics-consistency}

\subsubsection{Feature-attribution methods}\label{sec:feature-attribution}

To understand what the models learn, we examined gradient-based and
tree-based explanations. Grad-CAM was applied to ResNetVibCNN on the
vibration spectrogram only; TreeSHAP was applied to LightGBM on the fused
AE and vibration spectrograms. Both methods assign importance to input
frequency bins, but Grad-CAM does not provide an AE explanation.
Table~\ref{tbl:bandmass} summarises the distribution of importance mass
across physical frequency bands. For vibration, the bands separate
low-frequency wheel/spindle dynamics (\(<500\)~Hz), mid-frequency
structural response (500~Hz--2~kHz), the chatter/structural band where
wheel--workpiece interactions are prominent (2--15~kHz), and the
high-frequency vibration response (\(>15\)~kHz). For the vibration spectrogram, Grad-CAM and TreeSHAP
concentrate most importance in the chatter/structural and high-frequency
bands, with only small fractions below 2~kHz. AE is reported only with
modality-appropriate bands in Table~\ref{tab:ae-bands}; for the current
out-of-fold random forest 34.21\% of AE attribution mass lies above 1~MHz, outside
the historically reported nominal response band, and should be treated as a model
attribution rather than calibrated physical AE energy. For the recommended RF/dB-z model, TreeSHAP shows AE
importance concentrated above approximately 300~kHz and vibration
importance primarily between approximately 16 and 24~kHz; the AE importance above 1~MHz remains
sensor-response-limited (Figure~\ref{fig:rf-oof-shap}).

\begin{table}[htbp]
\centering
\caption{Importance mass by physical frequency
band.}\label{tbl:bandmass}
\begin{tabular}[]{@{}
  >{\raggedright\arraybackslash}p{(\linewidth - 10\tabcolsep) * \real{0.1515}}
  >{\raggedright\arraybackslash}p{(\linewidth - 10\tabcolsep) * \real{0.1212}}
  >{\raggedright\arraybackslash}p{(\linewidth - 10\tabcolsep) * \real{0.1667}}
  >{\raggedright\arraybackslash}p{(\linewidth - 10\tabcolsep) * \real{0.2273}}
  >{\raggedright\arraybackslash}p{(\linewidth - 10\tabcolsep) * \real{0.1667}}
  >{\raggedright\arraybackslash}p{(\linewidth - 10\tabcolsep) * \real{0.1667}}@{}}
\toprule
\begin{minipage}[b]{\linewidth}\raggedright
Modality
\end{minipage} & \begin{minipage}[b]{\linewidth}\raggedright
Method
\end{minipage} & \begin{minipage}[b]{\linewidth}\raggedright
\(<500\) Hz
\end{minipage} & \begin{minipage}[b]{\linewidth}\raggedright
500 Hz--2 kHz
\end{minipage} & \begin{minipage}[b]{\linewidth}\raggedright
2--15 kHz
\end{minipage} & \begin{minipage}[b]{\linewidth}\raggedright
\(>15\) kHz
\end{minipage} \\
\midrule
Vibration & Grad-CAM & 0.09\% & 2.26\% & 44.55\% & 53.10\% \\
Vibration & LightGBM global TreeSHAP & 3.55\% & 4.97\% & 13.33\% & 78.15\% \\
Vibration & RF OOF TreeSHAP & 0.21\% & 0.26\% & 10.89\% & 88.63\% \\
\bottomrule
\end{tabular}
\end{table}

\begin{figure}[htbp]
\centering
\includegraphics[width=183mm]{rf_dbz_oof_treeshap.png}
\caption{Primary interpretability analysis: out-of-fold TreeSHAP frequency
profiles for the reproducible current RF on AE-dB-z + Vib-dB-z. Curves show
sample-pooled attribution, equal-condition aggregation after fold-wise
normalisation, and the latter excluding Condition~7. Each held-out pass is
explained only by the RF trained without its condition. The analysis reproduces
the current 0.0210~\(\upmu\)m OOF MAE. Condition-balanced vibration attribution
above 15~kHz is 85.76\% (85.74\% without Condition~7), with a 16.2~kHz peak;
the AE profile peaks at 334~kHz.}\label{fig:rf-oof-shap}
\end{figure}

\textit{Note:} This shared-band table is restricted to vibration. AE
interpretation uses the modality-specific bands in Table~\ref{tab:ae-bands}.

\begin{table}[htbp]
\centering
\caption{Out-of-fold TreeSHAP importance mass for the reproducible current
random forest using AE-appropriate frequency bands.}\label{tab:ae-bands}
\begin{tabular}{@{}lc@{}}
\toprule
AE frequency band & Importance mass \\
\midrule
\(<100\)~kHz & 2.58\% \\
100--300~kHz & 6.59\% \\
300--800~kHz & 33.96\% \\
800~kHz--1~MHz & 22.66\% \\
\(>1\)~MHz & 34.21\% \\
\bottomrule
\end{tabular}
\end{table}

\subsubsection{Frequency-domain consistency with grinding mechanics}\label{sec:frequency-consistency}

Figure~\ref{fig:bandmass} visualises the same distribution. The
attribution distributions differ clearly between modalities: AE
explanations are almost entirely localised at high frequencies, whereas
vibration explanations are spread across the chatter and high-frequency
bands. The shared four-band taxonomy in Table~\ref{tbl:bandmass} is
retained only for visual comparison; AE-specific interpretation should
use separate bands such as \(<100\)~kHz, 100--300~kHz, 300--800~kHz,
800~kHz--1~MHz, and \(>1\)~MHz. Within vibration, Grad-CAM places
44.6\% of its mass in the 2--15~kHz chatter/structural band, while
RF out-of-fold TreeSHAP assigns 88.63\% to \(>15\)~kHz; the separate
LightGBM global TreeSHAP profile assigns 78.15\%. These differences cannot be
assigned to the explainer alone because model architecture, input
configuration, fitted predictor, and attribution aggregation also differ.

\begin{figure}
\centering
\includegraphics[width=183mm,height=\textheight,keepaspectratio]{results_interpretability_bandmass.png}
\caption{Vibration importance mass by physical frequency band for
ResNetVibCNN Grad-CAM, LightGBM global TreeSHAP, and RF out-of-fold
TreeSHAP. The differing model and input
configurations mean this is a descriptive comparison, not an isolated
comparison of explainers.}\label{fig:bandmass}
\end{figure}

\subsubsection{Physical interpretation of top features}\label{sec:top-features}

The dominant bins are compatible with grinding-related spectral activity,
but they do not identify a unique surface-generation mechanism. The
frequency-band ablation, sensor-response limitations, and absence of
archived no-cut spectra prevent a causal conclusion. Global LightGBM
TreeSHAP on vibration ranks bin 209 (\(20.9\) kHz)
first, accounting for 69.5\% of total importance mass. Grad-CAM peaks at
\(4.1\) kHz (bin 41), with the top ten bins spanning 3.7--12.0 kHz.
Global LightGBM TreeSHAP on AE peaks at \(1.22\) MHz (bin 182), which lies above the
historically reported 100~kHz--1~MHz nominal AE response band; this bin
should therefore be interpreted cautiously unless the sensor transfer
function above 1~MHz is independently verified. The top three AE bins
all lie above 840~kHz and together account for 55.3\% of AE importance
mass. A single dominant vibration bin contributes a large fraction of
importance. This is treated as upper-band predictive sensitivity; its
mechanical versus acquisition-chain origin cannot be resolved. Mounting,
sensor-response, preprocessing, or condition-specific artefacts cannot be
excluded. Fold-specific TreeSHAP was computed for the
reproducible current RF/dB-z pipeline using only each held-out condition's
predictions. The pooled out-of-fold profile peaks at 16.2~kHz for vibration
and 334~kHz for AE; vibration peaks remain mainly in the 16--24~kHz region,
whereas the dominant AE bin varies more strongly among folds
(Figure~\ref{fig:rf-oof-shap}). The current refits reproduce
an overall OOF MAE of 0.0210~\(\upmu\)m. The separate faster impurity-based
diagnostic in Supplementary Figure~\ref{fig:supp-shap-stability} provides a
qualitatively related stability check but is not SHAP. The AE attribution
should therefore be interpreted more cautiously. Additional
sensor-transfer-function and background-noise checks are needed to
exclude instrumentation artefacts. Because no-cut spectra and definitive
attribution stability are unavailable, the physical interpretation should
be treated as plausible but not validated.

Figure~\ref{fig:xai} illustrates model- and modality-specific explanations
from TreeSHAP (LightGBM on AE-dB + Vib-dB) and Grad-CAM
(ResNetVibCNN on Vib-dB). The TreeSHAP panel shows sparse AE importance
with labels converted from internal feature indices to MHz, while the
Grad-CAM panel shows a broader vibration-frequency profile. The profiles
share a broad separation between higher-frequency AE information and
structural-frequency vibration information, but their dominant bins differ
substantially by model and representation.

\begin{figure}
\centering
\includegraphics[width=183mm,height=\textheight,keepaspectratio]{shap_importance_LightGBMModel_ae_spec+vib_spec.png}
\caption{Explainability analysis. Left: TreeSHAP mean absolute
importance for the top AE frequency bins (LightGBM on
AE-dB + Vib-dB), labelled by physical frequency.
Right: Grad-CAM frequency
importance for ResNetVibCNN on Vib-dB. The panels illustrate modality- and model-specific attribution patterns and should not be interpreted as direct agreement between explainers. AE attributions above 1 MHz should be interpreted cautiously. The ablation tests are reported separately for their exact frequency ranges and do not make the attributed bands causally essential.}\label{fig:xai}
\end{figure}

\subsection{Uncertainty
quantification}\label{uncertainty-quantification}

\subsubsection{Coverage failure of MC-dropout epistemic-only intervals}\label{sec:aggregate-coverage}

MC-dropout was applied to ResNetVibCNN on Vib-dB under the
same 16-fold LOGO protocol used for accuracy, so every condition,
including Condition~7, appears as a held-out test set once under grouped
condition shift. Fifty
stochastic forward passes were run with dropout active and
batch-normalisation frozen in evaluation mode so that small test batches
did not shift the population statistics. The 95\% intervals are
\(\hat y \pm 1.96\,\sigma_{\text{MC}}\). Across all 319 held-out
LOGO samples the aggregate coverage is 50.2\%, the mean interval width
is \(0.0479~\upmu\)m, and the mean absolute error of the stochastic mean
prediction is \(0.0452~\upmu\)m. The exact same checkpoints give a
deterministic mean fold MAE of 0.0443~\(\upmu\)m and a stochastic-mean
fold MAE of 0.0452~\(\upmu\)m; their mean absolute prediction shift is
0.0023~\(\upmu\)m. The separate canonical 0.0268~\(\upmu\)m result is not a
same-checkpoint comparator. Supplementary Table~\ref{tab:supp-mcdropout-deterministic}
reports the fold-specific deterministic MAE, stochastic-mean MAE, and
mean absolute prediction shift. Activating dropout changes the point
predictor only slightly for this checkpoint set; the principal failure is
the severe interval undercoverage, so the output should be treated as an
uncertainty diagnostic rather than a calibrated predictive model.
Condition~7 has 0\% coverage, with a mean interval width of
\(0.041~\upmu\)m and a mean absolute error of \(0.247~\upmu\)m
(Supplementary Table~\ref{tab:supp-mcdropout-full-logo}). Condition~1
also has 0\% coverage (MAE \(0.044~\upmu\)m), so the undercoverage is
not confined to the extreme Condition~7 outlier; it affects any
held-out condition whose residual exceeds the narrow dropout variance.
Coverage is at or below 25\% for six of the 16 conditions, including
the hardest Conditions 7 and 8, reflecting a general inability of
dropout variance to scale with grouped condition shift. MC-dropout
therefore does not provide a predictive interval for this CNN's observed targets,
but it does carry ranking information: the Pearson correlation between
predicted standard deviation and absolute error is 0.48, the Spearman
correlation is 0.61, and the area under the ROC curve for detecting the
top 20\% highest-error samples is 0.82.

\begin{figure}
\centering
\includegraphics[width=183mm,height=\textheight,keepaspectratio]{mc_dropout_intervals_ResNetVibCNN_vib_spec.png}
\caption{Full 16-fold LOGO MC-dropout intervals for
ResNetVibCNN on Vib-dB. Blue band: nominal 95\%
MC-dropout epistemic interval; orange dots: observed \(R_a\). Samples are sorted by predicted
mean. Descriptive target coverage is 50.2\% across 319 held-out LOGO samples; the
mean interval width is 0.0479~\(\upmu\)m.}\label{fig:uncertainty}
\end{figure}

Figure~\ref{fig:uncertainty} shows the intervals sorted by predicted
mean. The intervals are far too narrow: many observed points lie outside
the band, especially in the difficult held-out conditions. The failure is
not restricted to Condition~7; coverage exceeds 75\% for five conditions
and reaches 75\% for one additional condition, but it is at or below
25\% for six conditions. This result
overturns the earlier 59-sample case study (Conditions 1, 2, and 6),
which reported 100\% coverage with a mean width of 0.073~\(\upmu\)m;
the apparent calibration of MC-dropout is therefore highly sensitive to
which conditions are included.

\subsubsection{Reliability by uncertainty bins}\label{sec:reliability-bins}

\begin{figure}
\centering
\includegraphics[width=183mm,height=\textheight,keepaspectratio]{mc_dropout_reliability.png}
\caption{Reliability diagram for full 16-fold LOGO MC-dropout on
ResNetVibCNN on Vib-dB. \textbf{a}, Empirical coverage per
uncertainty bin versus the nominal 95\% level. \textbf{b}, Mean absolute error
and nominal interval half-width per bin. Coverage remains below the nominal level in most bins; interval
half-width is too small to cover the observed error in high-error bins and some difficult conditions show systematic prediction bias.}\label{fig:reliability}
\end{figure}

Figure~\ref{fig:reliability} examines reliability by binning the 319
samples into five nearly equal-sized bins by predicted standard
deviation. The left panel shows empirical coverage for each bin. Most bins
are below the nominal 95\% level, so MC-dropout is generally
overconfident. The right panel shows that the mean nominal 95\% interval half-width is
much smaller than the mean absolute error in the highest-error bins,
confirming that the stochastic-forward variance underestimates prediction
error for the hardest samples.

\subsubsection{Sample-level calibration}\label{sec:sample-calibration}

\begin{figure}
\centering
\includegraphics[width=89mm,height=\textheight,keepaspectratio]{results_uncertainty_calibration.png}
\caption{Full 16-fold LOGO MC-dropout predicted standard deviation
versus observed absolute error for ResNetVibCNN on
Vib-dB. Pass-level Pearson $r=0.488$, Spearman $\rho=0.594$; the
empirical coverage is 50.2\%. Points above the diagonal are
outside one predicted standard deviation; the nominal 95\% boundary
would be error \(=1.96\sigma\). Many points lie above the one-standard-deviation
diagonal, especially at higher error. Green points are covered by the
nominal interval and red points are uncovered.}\label{fig:uncertainty-calibration}
\end{figure}

Figure~\ref{fig:uncertainty-calibration} disaggregates the reliability
assessment to individual samples. The pass-level Pearson correlation between
predicted uncertainty and absolute error is 0.488 (Spearman 0.594; AUROC for
detecting the top 20\% highest-error samples 0.818), indicating that
MC-dropout has some ranking ability---larger dropout variance tends to
accompany larger error---but the 95\% intervals are too narrow to cover
the observed error for roughly half of the samples. Calibration and
discrimination are therefore distinct: MC-dropout can partially rank
samples by predicted difficulty yet still fail as a calibrated
prediction interval. It should not be used as a standalone process-warning
signal for this CNN. Because passes are clustered within 16 conditions,
these pass-level quantities are descriptive. Across the 16 condition means,
the correlation between mean predicted standard deviation and mean absolute
error is Pearson $r=0.492$ (Spearman $\rho=0.732$). A condition-cluster
bootstrap gives wide 95\% intervals for the pass-level Pearson correlation
([0.202, 0.951]), Spearman correlation ([0.252, 0.817]), and AUROC
([0.471, 0.992]); Supplementary Table~\ref{tab:supp-clustered-uncertainty}
reports the corresponding leave-one-condition-out sensitivity.

Because MC-dropout is evaluated on a CNN, we also constructed diagnostic
intervals for a separate 15-condition RF uncertainty variant
(AE-dB-z + Vib-dB-z). The interval is computed as
\(\hat y \pm c \, \sigma_{\text{tree}}\), where
\(\sigma_{\text{tree}}\) is the standard deviation of the individual
tree predictions and the scaling factor \(c\) is calibrated using
out-of-bag (OOB) residuals from the training folds so that the test-fold
residuals do not influence interval width. The resulting 95\% intervals
achieve 80.3\% sample-weighted coverage across all 319 LOGO test samples
(Figure~\ref{fig:rf-conformal}), with a mean width of
\(0.070~\upmu\)m. This is below the nominal 95\% level, indicating that
OOB-based calibration is anti-conservative on this small dataset and
produces intervals that are too narrow. Coverage is uneven across
conditions: Condition~7 has 0\% coverage, whereas several well-behaved
conditions reach 100\% (Supplementary
Table~\ref{tab:rf-oob-coverage}). The undercoverage suggests that OOB
residuals underestimate held-out-condition residuals under grouped
distribution shift; future calibration should be condition-level, not
sample-level. A more fundamental limitation is that Condition~7 is
uncovered by every calibration strategy tested. Condition~7 has the
largest negative prediction bias (mean predicted $R_a$ about 0.217~\(\upmu\)m
versus mean measured 0.3540~\(\upmu\)m) and its observed $R_a$ lies above the
range seen in the training conditions, so the failure combines a
systematic directional bias with extrapolation outside the training-target
support. Condition~8, the next hardest condition, has the widest
calibrated interval (0.291~\(\upmu\)m) and 100\% coverage because its
observed $R_a$ remains inside the training-target range; Condition~7's
interval, although the second widest, is still too narrow because the
target lies outside that range. Tree variance alone cannot widen the
interval enough to cover this shift; future calibration must explicitly
model condition-level bias or reserve held-out conditions at the
parameter-space frontier.
We therefore compared the global OOB-scaled interval with three alternatives: a condition-aware factor calibrated from the three
nearest training conditions in normalised process-parameter space, a
split-conformal factor calibrated on one held-out condition, and a naive
residual interval using the 95th percentile of absolute OOB residuals
without tree-variance scaling. Table~\ref{tab:calibration-strategies}
shows that the global OOB-scaled interval actually gives the highest
mean condition-level coverage (80.3\%); the condition-aware and
split-conformal variants achieve 77.8\% and 72.4\%, respectively, and the
naive residual interval only 67.5\%. The split-conformal variant is
limited by having only one held-out calibration condition per fold, so it
cannot reliably calibrate condition-level tails. The interval uses the same
held-out LOGO test conditions as the accuracy evaluation, but its RF is fitted
on all 15 non-test conditions rather than the 14-condition current benchmark
split. It is therefore a separate uncertainty variant and a post-hoc
diagnostic, not an interval for the exact recommended RF or a full conformal
prediction guarantee.

\begin{table}[htbp]
\centering
\caption{Comparison of nominal 95\% interval strategies for a 15-condition
RF uncertainty variant. Coverage and width are averaged over the 16
LOGO test folds (mean per-condition coverage; the sample-weighted
aggregate for the global OOB strategy is the same to one decimal place,
80.3\%).}\label{tab:calibration-strategies}
\begin{tabular}{@{}lcc@{}}
\toprule
Strategy & Mean coverage & Mean width (\(\upmu\)m) \\
\midrule
Global OOB scaled & 80.3\% & 0.0704 \\
Condition-aware k-NN & 77.8\% & 0.0693 \\
Split conformal (1 calib. condition) & 72.4\% & 0.0810 \\
Naive residual 95\% & 67.5\% & 0.0264 \\
\bottomrule
\end{tabular}
\end{table}

\begin{figure}
\centering
\includegraphics[width=183mm,height=\textheight,keepaspectratio]{rf_conformal_intervals.png}
\caption{OOB-calibrated nominal 95\% intervals for a 15-condition RF
uncertainty variant (AE-dB-z + Vib-dB-z). Blue band: OOB-scaled interval; blue
dots: observed \(R_a\). Samples are sorted by predicted mean. Coverage is
80.3\% across all 319 held-out LOGO samples; mean width is
0.070~\(\upmu\)m.}
\label{fig:rf-conformal}
\end{figure}

\subsection{Condition-level failure analysis}\label{sec:conditions}

\subsubsection{Hard conditions and error-budget decomposition}\label{sec:hard-conditions}

Table~\ref{tbl:conditions} lists the grinding conditions with the
highest mean MAE across all model-configuration pairs. The ranking is
computed over 16 distinct conditions and all trained
model--configuration pairs. Condition~7 is a robust prediction failure:
its mean MAE of \(0.234~\upmu\)m is 2.2 times that of Condition~8 and
5.8 times that of Condition~10. Across the evaluated model--configuration
inventory, it accounts for 35.2\% of the summed condition-level mean MAEs;
this inventory-weighted quantity depends on the included configurations.
For the recommended RF alone, Condition~7 is discussed separately through
its held-out mean MAE and target-support diagnostic, while the top
three conditions (7, 8 and 10) together account for 57.4\%. Condition~7
also has the highest mean $R_a$ in the design, but its mean AE and
tri-axial vibration RMS values (0.437 and 1.591, respectively) are not
the highest (Supplementary Table~\ref{tab:supp-condition-diagnostics}),
ruling out a simple broadband RMS explanation. This does not exclude spectral, burst, thermal, force, coolant, or wheel-state mechanisms. By contrast, the median
condition MAE across all 16 conditions is only \(0.024~\upmu\)m, and
the remaining 11 conditions fall below \(0.030~\upmu\)m. For the canonical
recommended RF, descriptively excluding Condition~7 reduces mean fold MAE
from 0.0200 to 0.0118~\(\upmu\)m. Condition~7 remains part of every
principal result and should not be removed unless it
is proven to be invalid data, because difficult conditions are often the
most important cases in process monitoring.

\begin{table}[htbp]
\centering
\caption{Hardest grinding conditions (mean MAE across all
models).}\label{tbl:conditions}
\begin{tabular}[]{@{}ll@{}}
\toprule
Condition ID & Mean MAE (\(\upmu\)m) \\
\midrule
7 & 0.2337 \\
8 & 0.1073 \\
10 & 0.0402 \\
11 & 0.0343 \\
9 & 0.0300 \\
\bottomrule
\end{tabular}
\end{table}

Condition~7 is a consistent prediction outlier across models, but direct
process evidence for a distinct regime is not yet available. The failure
motivates spectral, force, power, or wheel-state diagnostics. Because workpiece speed and depth of
cut are not fully crossed in the fractional design, the Condition~7
interpretation is partly confounded with interaction effects that the
16-condition experiment cannot cleanly isolate. Until additional process
evidence is collected, Condition~7 should be treated as a difficult
prediction condition rather than as a confirmed process-regime outlier.
Conditions 8 and 10 are also elevated: Condition~8 has a mean MAE of
\(0.107~\upmu\)m, more than four times the median, and Condition~10 has
a mean MAE of \(0.040~\upmu\)m, still 1.7 times the median. These three
conditions are the only ones that lie clearly above the central cluster,
suggesting they represent distinct difficult grinding conditions rather
than random prediction noise.
For the recommended RF, absolute-residual Spearman correlation with run
order is 0.220, weaker than the corresponding correlations with
training-target support distance (0.429) and measured \(R_a\) (0.291);
therefore, the Condition~7 failure is not explained by a simple monotonic
run-order trend.

Figure~\ref{fig:condition-bars} shows the full condition ranking. The
spike is observed for the best random forest, ResNetVibCNN, and the best
LightGBM configuration, confirming that Condition~7 is a robust
prediction failure rather than an artefact of any single model. The
visual emphasises how disproportionately Condition~7 contributes to the
error budget---its bar is more than twice the height of the next largest and roughly an
order of magnitude above the median---so the practical priority is to identify
and monitor the few high-contribution conditions rather than to squeeze average MAE across
the remaining conditions. Because these condition-level differences have
direct implications for which model can be deployed, the next subsection
compares the deployment characteristics of the top candidates.

\begin{figure}
\centering
\includegraphics[width=183mm]{results_condition_error_bars.png}
\caption{Mean MAE across all model-configuration pairs for each
grinding condition, sorted from lowest to highest. Condition~7 is
highlighted as the largest outlier.}\label{fig:condition-bars}
\end{figure}

\begin{figure}
\centering
\includegraphics[width=183mm]{top_models_fold_mae.png}
\caption{Per-fold MAE for the top five model--configuration pairs. The
vertical dashed line marks condition 7, which produces the largest errors
for every leading model.}\label{fig:top-models-fold-mae}
\end{figure}

\begin{figure}
\centering
\includegraphics[width=183mm]{top_models_per_condition_mae.png}
\caption{Per-condition fold MAE heatmap for the top five
model--configuration pairs. Condition 7 stands out as the dominant source
of fold-level error.}\label{fig:top-models-per-condition}
\end{figure}

\subsection{Robustness and ablation
experiments}\label{sec:robustness}

\subsubsection{Cross-validation schemes}

Table~\ref{tab:cv-schemes} compares the separate 15-condition RF robustness
pipeline under
the available validation schemes. Standard 16-fold LOGO and the
workpiece-speed and grinding-depth grouped LOGO schemes give MAEs around
0.020--0.021~\(\upmu\)m. In contrast, LOGO grouped by wheel speed gives
0.0382~\(\upmu\)m, indicating that generalisation across wheel-speed
groups is materially harder for this model than generalisation across the
other two grouped process-parameter axes. This result should be
interpreted cautiously because each wheel-speed fold aggregates only four
design conditions, but the per-group results show that the 30~m/s group
(Conditions 5--8, including Condition~7) drives the degradation
(MAE 0.1088~\(\upmu\)m). The 25 and 40~m/s groups are both near
0.009~\(\upmu\)m, and the 35~m/s group reaches 0.0260~\(\upmu\)m.
Thus, the grouped result reflects both wheel-speed transfer difficulty and
the known high-error Condition~7 regime; it is not evidence of a uniform
wheel-speed effect across all four levels.
Random 5-fold CV and random 80/20 splits give
much lower MAEs (0.0053~\(\upmu\)m), which is below the
manufacturer-quoted repeatability and illustrates the optimism that
arises when highly similar passes from the same grinding condition appear in both training and test sets. Such random-split MAEs should not be interpreted
as physically meaningful predictive performance. The grouped designs are
therefore essential for an honest estimate of generalisation.

\begin{table}[htbp]
\centering
\caption{Validation-scheme comparison for the 15-condition RF robustness
pipeline (AE-dB-z + Vib-dB-z). For each grouped scheme, the complete indicated
group is withheld and all remaining groups are used for fitting. This is not
the 14-condition current benchmark RF.}\label{tab:cv-schemes}
\begin{tabular}{@{}lrc@{}}
\toprule
Scheme & Folds & MAE (\(\upmu\)m) \\
\midrule
LOGO (16 conditions) & 16 & 0.0202 \\
LOGO by workpiece speed & 4 & 0.0210 \\
LOGO by grinding depth & 4 & 0.0204 \\
LOGO by wheel speed & 4 & 0.0382 \\
Random 5-fold CV (mean of 5 seeds) & 5 & 0.0053 \\
Random 80/20 split (mean of 10 seeds) & 10 & 0.0053 \\
\bottomrule
\end{tabular}
\end{table}

\begin{table}[htbp]
\centering
\caption{Per-group errors for wheel-speed grouped validation of the
recommended random forest. Each group is withheld as four complete design
conditions.}\label{tab:wheel-speed-groups}
\begin{tabular}{@{}lcc@{}}
\toprule
Wheel speed (m/s) & Held-out conditions & MAE (\(\upmu\)m) \\
\midrule
25 & 1--4 & 0.0088 \\
30 & 5--8 & 0.1088 \\
35 & 9--12 & 0.0260 \\
40 & 13--16 & 0.0086 \\
\bottomrule
\end{tabular}
\end{table}

The 0.0200~\(\upmu\)m value in Table~\ref{tbl:top10} is the canonical
top-ranking LOGO summary. The strict regenerated normalisation ablation
uses the recorded seed-42 LOGO validation exclusions, current cache hashes,
and the documented RF settings, but gives 0.0210~\(\upmu\)m for dB-z
inputs and 0.0213~\(\upmu\)m for raw dB inputs. These values and their
experiment identifier are reported in Supplementary
Table~\ref{tab:supp-psnorm-ablation}. The 0.0010~\(\upmu\)m difference
from the canonical artifact means that this ablation is a reproducible
current-cache sensitivity analysis, not an exact reproduction of the
earlier canonical prediction artifact; it remains below the stated
practical metrology scale.

\subsubsection{Pass-position sensitivity}

The cached descriptor averages local spectrogram maps over an entire pass,
which removes temporal ordering. To test whether this averaging hides a
large position effect, the recommended RF was retrained on the early,
middle, or late third of the available local-map axis for every pass,
using the same 16 seed-42 LOGO folds and per-sample dB-z normalisation.
The mean MAEs were 0.0213, 0.0204, and 0.0198~\(\upmu\)m, respectively
(Supplementary Table~\ref{tab:supp-pass-position}). Concatenating the three
ordered third descriptors before fitting gives 0.0195~\(\upmu\)m
(fold SD 0.0338~\(\upmu\)m). This small numerical difference is below the
practical metrology scale used for top-model ranking and does not validate a
streaming system: all variants are retrospective partitions of archived pass
data, and the cache does not retain the information needed to compare fixed
physical-distance windows.

For context, simple LOGO baselines that use only the training targets are
substantially weaker: the training-target mean gives 0.0453~\(\upmu\)m MAE,
the training-target median gives 0.0340~\(\upmu\)m, and the nearest
training condition in standardised process-parameter space gives
0.0591~\(\upmu\)m. These baselines confirm that the recommended RF improves
on trivial target-distribution and parameter-neighbour predictors, while
still failing on the high-roughness frontier.

\subsubsection{Sensor-modality ablation}

Table~\ref{tab:modality-ablation} reports an exact 14-condition current-pipeline
ablation when only AE, only vibration, or both modalities are used. Vibration
alone achieves 0.0217~\(\upmu\)m, close to the fused configuration
(0.0210~\(\upmu\)m), whereas AE alone degrades to 0.0282~\(\upmu\)m. Under
the current sensor channels, preprocessing, and unequal window lengths,
vibration therefore carries most of the predictive information. The
0.0008~\(\upmu\)m AE-fusion difference is below the stated practical
metrology scale and is not evidence of a practically meaningful average
accuracy gain.

\begin{table}[htbp]
\centering
\caption{Sensor-modality ablation for the current 14-condition dB-z RF under
16-fold LOGO. The canonical test and validation conditions are excluded in
every fold.}\label{tab:modality-ablation}
\begin{tabular}{@{}lc@{}}
\toprule
Input & MAE (\(\upmu\)m) \\
\midrule
AE only & 0.0282 \\
Vib only & 0.0217 \\
AE + Vib & 0.0210 \\
\bottomrule
\end{tabular}
\end{table}

\subsubsection{Frequency-band ablation}

Table~\ref{tab:frequency-ablation} reports the original 15-condition robustness-pipeline evaluation-time
masking experiment, in which selected standardized features are replaced by
their training-fold mean in an already trained RF. Removing the AE content
above 1~MHz or the vibration 2--15~kHz band individually changes MAE by
less than 0.0001~\(\upmu\)m; removing both gives a small numerical increase
to 0.0208~\(\upmu\)m. Supplementary
Table~\ref{tab:supp-retrained-band-ablation} then retrains the RF after
setting the selected dB-z features to zero before both fitting and testing.
Removing vibration above 15~kHz raises MAE from 0.0210 to
0.0285~\(\upmu\)m; the mean paired increase is 0.00755~\(\upmu\)m with a
descriptive condition-bootstrap 95\% interval of [0.00149, 0.01532]. Removing
only the narrower 16--24~kHz region gives 0.0234~\(\upmu\)m, but its paired
interval includes zero. Thus, the AE >1~MHz and vibration 2--15~kHz bands
are redundant under these RF experiments, whereas the broader vibration
region above 15~kHz contributes predictive information when removed in full.
Because those bins were zeroed after the original full-spectrum dB-z
normalisation, retained bins could still encode the removed band's contribution
to the sample mean and variance. A stricter sensitivity deletes the vibration
band from the dB cache before recomputing each sample/channel z-score
(Supplementary Table~\ref{tab:supp-prez-band-ablation}). The recomputed full
baseline gives 0.0209~\(\upmu\)m; removing >15~kHz gives 0.0255~\(\upmu\)m
(paired increase 0.00455, descriptive 95\% interval [0.00063, 0.01066]), and
removing 16--24~kHz gives 0.0254~\(\upmu\)m (increase 0.00447, interval
[-0.00035, 0.01050]). Thus, direct upper-band dB-z features are predictive,
and the broader >15~kHz sensitivity persists when indirect normalisation
information is removed.
These ablations do not establish that any band is physically causal. In
particular, the >1~MHz AE result remains a model-response finding rather than
a calibrated physical-energy claim. The (>15)~kHz vibration region also
lies close to the 25.6~kHz Nyquist limit. Because the accelerometer response,
analogue anti-alias filter, and mounting transfer function were not archived,
its predictive value cannot be assigned to a mechanical mechanism rather
than acquisition-chain roll-off or resonance.

\begin{table}[htbp]
\centering
\caption{Evaluation-time frequency-band masking for the 15-condition robustness RF under 16-fold LOGO. Only the test condition is excluded; this is not an exact ablation of the 14-condition current benchmark RF.}\label{tab:frequency-ablation}
\begin{tabular}{@{}lc@{}}
\toprule
Input & MAE ($\upmu$m) \\
\midrule
Full AE + Vib & 0.0202 \\
AE \textgreater 1~MHz removed & 0.0203 \\
Vib 2--15~kHz removed & 0.0202 \\
Both removed & 0.0208 \\
\bottomrule
\end{tabular}
\end{table}


\subsection{Deployment
characteristics}\label{deployment-characteristics}

\subsubsection{Latency--accuracy frontier}\label{sec:latency-accuracy}

The present deployment analysis is limited to feature-precomputed
model-inference timing; full online latency from raw signal to prediction
remains future work. Table~\ref{tbl:latency} summarises single-sample CPU
inference latency after feature extraction and
checkpoint size for representative top model--configuration pairs. The
timings were measured on a single logical core of an Intel Xeon Silver
4514Y (2.0~GHz base, 3.4~GHz max turbo, Ubuntu 22.04 with kernel
5.15.0-179, `powersave' governor, PyTorch 2.5.1, scikit-learn 1.6.1) with
all threading libraries pinned to one thread, batch size one, and 100 warm-up repetitions excluded; garbage collection was disabled during each timing block, and the benchmark was repeated in five independent blocks of 200 runs each. The results quantify model inference only and do not include signal acquisition, STFT/dB-z/log-mel computation, feature scaling, communication overhead, or edge-hardware scheduling. These
software versions are those
of the latency-measurement environment; the original LOGO
training/evaluation runs were performed under earlier package versions
and may differ in minor dependency details. The `powersave' governor can
increase latency variability, so the reported p95 values are timing
diagnostics under the measurement environment and should not be
interpreted as hard real-time bounds. They should be repeated under a fixed
performance governor for a hard real-time assessment. Each row pools timing
over the same 16 canonical repeat-0 outer-fold checkpoints that produce its
reported 16-fold MAE. Every checkpoint is timed in five independent blocks of
200 calls using the first held-out sample from its corresponding condition,
giving 16,000 calls per configuration. Checkpoint paths, complete SHA-256
digests, and fold-specific timing summaries are provided in
\texttt{edge\_latency\_benchmark\_per\_fold.csv}. Random forest on dB-z has
a pooled median latency of 8.0 ms, p95 latency of 13.7 ms, and median
checkpoint size of 2.92 MB. The RF on AE-dB + Vib-dB has a similar pooled
median of 8.0 ms and a tighter p95 of 8.2 ms.

The CNNs are the fastest options: ResNetVibCNN on Vib-dB
has a pooled median latency of 1.3 ms and ResNetAECNN on AE-dB of
2.3 ms, and their median checkpoints are the smallest (0.35 MB). Their
point-estimate MAEs are \(0.0068\)--\(0.0100~\upmu\)m above the historical
best RF value; most corresponding pairwise comparisons are not significant
after multiplicity correction. LightGBM on \texttt{Vib-dB-z}
sits between the two: MAE \(0.0250~\upmu\)m and pooled median latency 4.9 ms,
with a p95 of 5.0 ms; it remains less accurate than the best RF while
offering a much smaller checkpoint. Bilinear fusion on
AE-dB + Vib-dB + phys + PP is the largest deep-learning
entry (171k parameters, 4.0 ms median latency, 4.0 ms p95). This is its
overall best configuration, distinct from the time-domain-feature
variant in Table~\ref{tbl:representations}; even so, its accuracy
(\(0.0317~\upmu\)m) is no better than the much simpler CNNs.

\begin{table}[htbp]
\centering
\caption{CPU feature-precomputed model-inference latency and checkpoint
storage footprint. Each row combines the 16-fold benchmark MAE with timing
pooled over the same 16 historical canonical repeat-0 fold checkpoints. The
timed dB-z RF row therefore corresponds to the historical 0.0200~\(\upmu\)m
artifact, not the reproducible current 0.0210~\(\upmu\)m checkpoint set; the
current set has the same architecture and feature dimensions but was not timed.
Each checkpoint
uses five independent blocks of 200 runs, with warm-up excluded and garbage
collection disabled; p50 is the pooled median. Checkpoint size is the median
serialized size across the 16 checkpoints, not peak runtime memory. Tree
ensembles do not have gradient-trained neural weights, but they store
learned split thresholds, feature indices, and leaf predictions.}
\label{tbl:latency}
\scriptsize
\resizebox{\linewidth}{!}{%
\begin{tabular}{@{}llccccccc@{}}
\toprule
Model & Config & Benchmark MAE (\(\upmu\)m) & p50 (ms) & p90 (ms) & p95 (ms) & p99 (ms) & Max (ms) & Median checkpoint (MB) \\
\midrule
Random forest & \texttt{AE-dB-z + Vib-dB-z} & 0.0200 & 8.0 & 8.2 & 13.7 & 14.0 & 15.3 & 2.92 \\
Random forest & AE-dB + Vib-dB & 0.0210 & 8.0 & 8.1 & 8.2 & 8.3 & 10.8 & 2.74 \\
LightGBM & \texttt{Vib-dB-z} & 0.0250 & 4.9 & 5.0 & 5.0 & 5.0 & 7.7 & 0.70 \\
ResNetVibCNN & Vib-dB & 0.0268 & 1.3 & 1.3 & 1.3 & 1.4 & 4.1 & 0.35 \\
ResNetAECNN & AE-dB & 0.0300 & 2.3 & 2.3 & 2.3 & 2.9 & 5.2 & 0.35 \\
Bilinear fusion & AE-dB + Vib-dB + phys + PP & 0.0317 & 4.0 & 4.0 & 4.0 & 4.2 & 6.9 & 0.72 \\
\bottomrule
\end{tabular}
}
\end{table}

Supplementary Table~\ref{tab:supp-runtime-memory} reports a separate
fresh-process resident-memory diagnostic. After cached inputs are loaded,
model loading and one warm-up inference add 5--7~MB for the tree models,
14--15~MB for the single-modality CNNs, and 20~MB for bilinear fusion.
This is an incremental resident-memory measurement under the timing
environment, not peak deployment memory or an end-to-end edge-hardware
requirement.

The latency--accuracy trade-off is visualised in Figure~\ref{fig:pareto}.
Under the observed feature-precomputed inference comparison, the evaluated
tree-based configurations occupy the lowest-error region of the frontier: no
CNN reaches their low error, while LightGBM, bilinear fusion, and both CNNs
achieve approximately sub-5-ms pooled median model inference. These
configurations therefore occupy the faster-but-less-accurate region of the
plot. After cached features are available, every listed configuration exceeds
50 model calls per second by pooled median latency; this is not a 50~Hz
process-monitoring update rate. ResNetVibCNN provides the lowest median and
p95 latency. The reported cache
averages 2,910 local windows over the recorded pass, so its observation
horizon is the full pass (approximately 14.6~s in the raw AE record), not
110~ms. A streaming system with shorter or overlapping aggregation windows
would require a separate end-to-end evaluation. When the priority is minimising
point-estimate prediction error under the current evaluation, the random
forest is the preferred candidate. The deployment analysis is limited to
feature-precomputed model-inference timing; full online latency remains
future work.

\begin{figure}
\centering
\includegraphics[width=183mm]{latency_vs_accuracy.png}
\caption{Feature-precomputed latency versus accuracy for representative
models. Points use median latency; rightward horizontal whiskers extend to
p95 latency from Table~\ref{tbl:latency}. The y-axis shows mean LOGO MAE
without symmetric fold-standard-deviation bars because MAE is nonnegative
and the fold distribution is tail-sensitive.}\label{fig:pareto}
\end{figure}

\subsubsection{Footprint and maintainability}\label{sec:footprint-maintainability}

Figure~\ref{fig:deployment-tradeoffs} places the same models in a
three-way accuracy--latency--footprint space. Bubble size encodes
checkpoint size and colour distinguishes model configurations. The two random
forest configurations remain the preferred point-estimate candidates
under feature-precomputed model-inference timing: they combine the lowest
error with
pooled medians of approximately 8 ms, p95 values of 8.2--13.7 ms, and the
largest median checkpoints (2.7--2.9 MB). The
tree ensembles store learned split thresholds, feature indices, and leaf
predictions rather than gradient-trained weights. Their stored structure permits inspection of split rules and feature attributions, whereas retraining effort and engineer auditability were not measured. LightGBM lies in the middle of the latency--error
plane with a moderate median checkpoint (0.70 MB) and a pooled p95 of 5.0 ms.
The deep-learning CNNs cluster in the fast, small-footprint
lower-left region with \(\sim\)86k parameters, indicating that their
advantages are deployment size and speed, not accuracy under this LOGO
protocol. These trade-offs frame the practical recommendations discussed
in Section~\ref{discussion}.

\begin{figure}
\centering
\includegraphics[width=183mm]{results_deployment_tradeoffs.png}
\caption{Deployment trade-offs for representative models using the exact
Table~\ref{tbl:latency} median latencies. Bubble size is proportional to
checkpoint size and colour distinguishes model configurations; numbered markers identify
models in the legend below the plot. Auditable tree ensembles occupy the
accuracy-efficient region of the frontier.}\label{fig:deployment-tradeoffs}
\end{figure}
